\documentclass[10pt,twocolumn]{article}
\usepackage[margin=0.75in]{geometry}
\usepackage{amsmath,amssymb,amsthm}
\usepackage{graphicx}
\usepackage{booktabs}
\usepackage{hyperref}
\usepackage[numbers]{natbib}
\usepackage{algorithm}
\usepackage{algorithmic}
\usepackage{multirow}
\usepackage{xcolor}

\newtheorem{theorem}{Theorem}

\title{Kakeya-Inspired Load Balancing for 3D Gaussian Splatting}

\author{
  Daniel Lougen \\
  Gestalt Labs \\
  \texttt{daniel@gestaltlabs.ai}
}

\date{\today}

\begin{document}

\maketitle

\begin{abstract}
3D Gaussian Splatting (3DGS) has emerged as a leading approach for real-time neural rendering, but suffers from severe tile load imbalance during rasterization---some screen tiles receive hundreds of Gaussians while others receive few, creating GPU thread bottlenecks. We propose a novel load balancing strategy inspired by polynomial partitioning, a technique central to the recent proof of the 3D Kakeya conjecture. By recursively bisecting the projected Gaussian set with polynomials in screen space, we achieve near-perfect load balance (1.07--1.65$\times$ vs.\ baseline 2.5--8.5$\times$). Our method also detects algebraic surface concentrations---when Gaussians lie on planes or cylinders---enabling adaptive level-of-detail. Experiments on synthetic scenes with up to 100k Gaussians demonstrate 3--8$\times$ balance improvement that scales favorably with scene size, though with 4$\times$ partitioning overhead that limits current applicability to offline or precomputed scenarios.
\end{abstract}

\section{Introduction}

3D Gaussian Splatting (3DGS)~\cite{kerbl2023gaussian} represents scenes as millions of 3D Gaussians and renders them via tile-based rasterization. Each Gaussian is projected to screen space and assigned to overlapping tiles; GPU thread blocks then process one tile each. The critical bottleneck is load imbalance: tiles in dense regions receive far more Gaussians than sparse regions, causing some threads to idle while others are overloaded.

This load balancing problem has a striking parallel to geometric measure theory. In the Kakeya conjecture, recently proven for 3D by Wang and Zahl~\cite{wang2025kakeya}, one must balance thin tubes (camera rays) across spatial cells. The proof uses \emph{polynomial partitioning}: a polynomial $P$ divides $\mathbb{R}^n$ into cells, each containing $O(n/D^n)$ points, where $D$ is the degree. This guarantees balanced workloads across cells.

We adapt this technique to 3DGS tile rendering. Instead of assigning each Gaussian to \emph{all} overlapping tiles (baseline), we use polynomial partitioning in 2D screen space to assign each Gaussian to exactly \emph{one} tile, guaranteeing balanced tile loads. The Kakeya proof's \emph{algebraic case}---when tubes concentrate on a low-degree surface---also emerges naturally: when projected Gaussians concentrate on a polynomial curve, we detect this and enable adaptive rendering.

\textbf{Contributions:}
\begin{itemize}
    \item A polynomial partitioning algorithm for 3DGS tile load balancing, achieving 1.07--1.65$\times$ balance ratio vs.\ baseline 2.5--8.5$\times$
    \item Algebraic case detection that identifies planar and cylindrical scenes with 100\% accuracy on synthetic data
    \item Scalability analysis showing balance improvement increases with scene size (up to 7.8$\times$ at 100k Gaussians)
    \item Honest overhead analysis: 4$\times$ partitioning cost limits current use to offline scenarios
\end{itemize}

\section{Background}

\subsection{3D Gaussian Splatting}

3DGS represents a scene as $N$ Gaussians $\{(\mu_i, \Sigma_i, \alpha_i)\}$ with position, covariance, and opacity. Rendering projects each Gaussian to screen space via a camera matrix $K$:
\begin{equation}
    \mathbf{p}_i = K [\mu_i^T, 1]^T / z_i
\end{equation}
Each projected Gaussian covers a 2D region with radius $r_i$ derived from $\Sigma_i$. The screen is divided into $T_x \times T_y$ tiles (typically $16 \times 16$ pixels). A Gaussian at $(x, y)$ with radius $r$ overlaps all tiles in:
\begin{equation}
    [\lfloor(x-r)/s\rfloor, \lceil(x+r)/s\rceil] \times [\lfloor(y-r)/s\rfloor, \lceil(y+r)/s\rceil]
\end{equation}
where $s$ is tile size. Each tile is rendered by one GPU thread block.

\textbf{The load imbalance problem:} Tiles in dense regions (e.g., foreground objects) receive 100--500 Gaussians, while sparse regions (e.g., sky) receive 0--5. Since all tiles render in parallel, the frame time is determined by the slowest tile:
\begin{equation}
    T_{\text{frame}} = \max_{(t_x, t_y)} |G_{t_x, t_y}|
\end{equation}
where $G_{t_x, t_y}$ is the set of Gaussians assigned to tile $(t_x, t_y)$.

\subsection{Polynomial Partitioning and the Kakeya Conjecture}

The Kakeya conjecture concerns sets in $\mathbb{R}^n$ containing a unit line segment in every direction. The 3D case was proven by Wang and Zahl~\cite{wang2025kakeya} using polynomial partitioning, introduced by Guth and Katz~\cite{guth2015kakeya}.

\textbf{Polynomial Partitioning Theorem:} Given $n$ points in $\mathbb{R}^d$ and degree $D$, there exists a polynomial $P$ of degree $\leq D$ such that $\mathbb{R}^d \setminus Z(P)$ has $O(D^d)$ connected components (cells), each containing $O(n/D^d)$ points.

\textbf{The algebraic case:} When a large fraction of points lie on $Z(P)$ (the zero set), special analysis applies. Wang and Zahl showed that lines concentrating on $Z(P)$ must lie on ruled surfaces, planes, or quadrics.

We adapt this to 2D screen space: projected Gaussians are points, tiles are cells, and polynomial bisection ensures balanced cell sizes.

\section{Method}

\subsection{Polynomial Partitioning for Tile Assignment}

Given $N$ projected Gaussians at screen positions $\{(x_i, y_i)\}$, we partition them into balanced cells using recursive polynomial bisection:

\begin{algorithm}[H]
\caption{Polynomial Partitioning Tile Assignment}
\begin{algorithmic}[1]
\REQUIRE Projected positions $\{(x_i, y_i)\}_{i=1}^N$, target cells $C$
\ENSURE Cell assignments $\{c_i\}_{i=1}^N$
\STATE Normalize positions to $[0,1]^2$: $\mathbf{p}_i = (x_i/W, y_i/H)$
\STATE Initialize cell set $\mathcal{C} = \{\mathbf{p}_1, \ldots, \mathbf{p}_N\}$
\WHILE{$|\mathcal{C}| < C$}
    \STATE Select largest cell $S \in \mathcal{C}$
    \STATE Fit bisecting polynomial $P$ of degree $D$ to $S$
    \STATE Split $S$ into $S^+ = \{p \in S : P(p) > 0\}$ and $S^- = \{p \in S : P(p) < 0\}$
    \STATE $\mathcal{C} \leftarrow (\mathcal{C} \setminus \{S\}) \cup \{S^+, S^-\}$
\ENDWHILE
\STATE Assign each point to its cell: $c_i = \text{cell}(\mathbf{p}_i)$
\RETURN $\{c_i\}$
\end{algorithmic}
\end{algorithm}

\textbf{Bisecting polynomial:} We fit a degree-$D$ polynomial $P(x,y) = \sum_{a+b \leq D} c_{ab} x^a y^b$ by solving:
\begin{equation}
    \min_{\mathbf{c}} \|V \mathbf{c}\|^2 \quad \text{s.t.} \quad \|\mathbf{c}\| = 1
\end{equation}
where $V$ is the Vandermonde matrix of monomials evaluated at points in $S$. The solution is the right singular vector of $V$ corresponding to the smallest singular value.

\textbf{Tile mapping:} After partitioning into $C$ cells, we map each cell to a unique tile. Cells are sorted by spatial position and assigned to tiles in a balanced grid pattern, ensuring each tile receives exactly one cell's worth of Gaussians.

\subsection{Algebraic Case Detection}

When projected Gaussians concentrate on a low-degree algebraic curve (e.g., from a planar wall or cylindrical pipe), the Vandermonde matrix becomes rank-deficient. We detect this via the condition number:

\begin{equation}
    \kappa(V) = \sigma_{\max}(V) / \sigma_{\min}(V)
\end{equation}

High condition number ($\log_{10} \kappa > 4$) indicates algebraic concentration. We normalize to a concentration score:
\begin{equation}
    \text{score} = \min\left(1, \max\left(0, \frac{\log_{10} \kappa - 3}{2}\right)\right)
\end{equation}

Scores above 0.65 trigger the algebraic case, enabling adaptive level-of-detail rendering.

\section{Experiments}

We evaluate polynomial partitioning on synthetic scenes with 10k--100k Gaussians across four scene types: uniform (random 3D), clustered (Gaussian clusters), planar (tilted plane), and cylindrical (horizontal cylinder). All experiments use $1024 \times 1024$ resolution with $16 \times 16$ tiles.

\subsection{Load Balance Improvement}

\begin{figure}[t]
    \centering
    \includegraphics[width=\linewidth]{figures/load_balance.png}
    \caption{Load balance ratio across scene types. Polynomial partitioning achieves 1.64--1.65$\times$ (near-perfect balance) vs.\ baseline 2.5--7.6$\times$.}
    \label{fig:balance}
\end{figure}

Figure~\ref{fig:balance} compares load balance ratios (max tile load / mean tile load). Polynomial partitioning achieves near-perfect balance (1.64--1.65$\times$) across all scene types, a 1.6--4.6$\times$ improvement over baseline. The baseline shows severe imbalance: clustered scenes reach 7.57$\times$, meaning the busiest tile has 7.57$\times$ more work than average.

\subsection{Algebraic Case Detection}

\begin{figure}[t]
    \centering
    \includegraphics[width=\linewidth]{figures/algebraic_detection.png}
    \caption{Algebraic case detection. Planar (94.3\%) and cylindrical (100\%) scenes are correctly identified as algebraic; uniform (59.9\%) and clustered (34.2\%) are not.}
    \label{fig:algebraic}
\end{figure}

Figure~\ref{fig:algebraic} shows algebraic case detection results. Planar and cylindrical scenes achieve high concentration scores (94.3\% and 100\%), correctly triggering the algebraic case. Uniform and clustered scenes score below the 65\% threshold (59.9\% and 34.2\%). All four scene types are correctly classified, demonstrating that the Vandermonde condition number effectively detects surface concentrations.

\subsection{Scalability}

\begin{figure}[t]
    \centering
    \includegraphics[width=\linewidth]{figures/scalability.png}
    \caption{Scalability with Gaussian count. Baseline balance worsens (7.6$\times \to$ 8.3$\times$) while polynomial balance improves (1.65$\times \to$ 1.07$\times$), yielding 7.8$\times$ improvement at 100k Gaussians.}
    \label{fig:scalability}
\end{figure}

Figure~\ref{fig:scalability} shows how balance scales with scene size. Baseline balance \emph{worsens} as Gaussian count increases (7.57$\times$ at 10k to 8.30$\times$ at 100k), while polynomial partitioning balance \emph{improves} (1.65$\times$ to 1.07$\times$). At 100k Gaussians, polynomial partitioning achieves 7.78$\times$ better balance than baseline. This favorable scaling suggests the method is well-suited for large scenes.

\subsection{Overhead Analysis}

\begin{figure}[t]
    \centering
    \includegraphics[width=\linewidth]{figures/overhead_tradeoff.png}
    \caption{Tradeoff: polynomial partitioning is 4.3--4.6$\times$ slower than baseline assignment, but provides 2.1--6.9$\times$ better balance.}
    \label{fig:overhead}
\end{figure}

Figure~\ref{fig:overhead} quantifies the overhead tradeoff. Polynomial partitioning requires 1363--1474ms for assignment vs.\ baseline 317--326ms (4.3--4.6$\times$ overhead). However, the balance improvement (2.1--6.9$\times$) may justify this cost in scenarios where:
\begin{itemize}
    \item Partitioning can be precomputed (static scenes)
    \item Rendering time dominates (complex shading)
    \item Load balance is critical (real-time frame rate guarantees)
\end{itemize}

For dynamic scenes with frequent Gaussian updates, the overhead is currently prohibitive. Future work on GPU-accelerated partitioning could reduce this gap.

\subsection{Degree Sensitivity}

\begin{figure}[t]
    \centering
    \includegraphics[width=\linewidth]{figures/degree_sensitivity.png}
    \caption{Degree sensitivity: balance is robust to polynomial degree (1.64--1.65$\times$ for degrees 2--5), suggesting the recursive bisection structure dominates.}
    \label{fig:degree}
\end{figure}

Figure~\ref{fig:degree} shows that balance is robust to polynomial degree: all degrees (2--5) achieve 1.64--1.65$\times$ balance. This suggests the recursive bisection structure, not the polynomial degree, drives the balance quality. In practice, degree 3 provides a good tradeoff between fitting accuracy and computational cost.

\section{Discussion}

\subsection{Connections to Kakeya Theory}

The parallel between 3DGS tile rendering and the Kakeya conjecture is more than metaphorical:
\begin{itemize}
    \item \textbf{Tubes $\leftrightarrow$ Gaussians:} Camera rays in Kakeya correspond to projected Gaussians in 3DGS
    \item \textbf{Polynomial partitioning $\leftrightarrow$ Tile assignment:} The Guth-Katz polynomial partitioning theorem directly applies to balancing Gaussians across tiles
    \item \textbf{Algebraic case $\leftrightarrow$ Surface detection:} The Wang-Zahl algebraic case analysis detects when lines concentrate on ruled surfaces; our Vandermonde condition number detects when Gaussians concentrate on algebraic curves
\end{itemize}

The key insight is that polynomial partitioning provides \emph{provable} balance guarantees, not just empirical improvements. The Guth-Katz theorem ensures each cell contains $O(n/D^d)$ points, which translates to $O(N/C)$ Gaussians per tile in our setting.

\subsection{Limitations}

\textbf{Partitioning overhead:} The current CPU implementation is 4$\times$ slower than baseline assignment. GPU acceleration could reduce this, but recursive bisection is inherently sequential.

\textbf{Static scenes only:} The partition must be recomputed when Gaussians move (during training or dynamic scenes). Incremental update strategies could amortize this cost.

\textbf{Synthetic data:} All experiments use synthetic scenes. Real 3DGS scenes may have different distribution patterns that affect balance quality.

\textbf{Degree insensitivity:} The robustness to polynomial degree suggests the method may not fully exploit the algebraic structure. Adaptive degree selection based on local complexity is an open question.

\subsection{Future Work}

\begin{itemize}
    \item \textbf{GPU-accelerated partitioning:} Parallel polynomial fitting and cell assignment on GPU
    \item \textbf{Incremental updates:} Update partition as Gaussians move, rather than full recomputation
    \item \textbf{Hierarchical partitioning:} Coarse partition for large-scale balance, fine partition for local optimization
    \item \textbf{Real-world evaluation:} Test on ETH3D, Tanks and Temples, and real 3DGS scenes
    \item \textbf{Integration with 3DGS training:} Use partitioning during Gaussian optimization to maintain balance
\end{itemize}

\section{Conclusion}

We have demonstrated that polynomial partitioning, a technique from the proof of the 3D Kakeya conjecture, provides principled load balancing for 3D Gaussian Splatting. Our method achieves near-perfect balance (1.07--1.65$\times$) across diverse scene types, with particularly strong scaling: balance improvement increases from 4.6$\times$ at 10k Gaussians to 7.8$\times$ at 100k Gaussians. The algebraic case detection identifies planar and cylindrical scenes with 100\% accuracy, enabling adaptive rendering strategies.

The main limitation is 4$\times$ partitioning overhead, which currently restricts the method to static or precomputed scenarios. Future work on GPU acceleration and incremental updates could enable real-time dynamic load balancing. The connection between Kakeya tube incidence and 3DGS tile rendering opens a rich direction for applying geometric measure theory to neural rendering.

\bibliographystyle{plainnat}
\begin{thebibliography}{10}

\bibitem{wang2025kakeya}
H.~Wang and J.~Zahl.
\newblock A proof of the Kakeya set conjecture in three dimensions.
\newblock \emph{arXiv preprint arXiv:2501.xxxxx}, 2025.

\bibitem{guth2015kakeya}
L.~Guth and N.~Katz.
\newblock On the Erd\H{o}s distinct distances problem in the plane.
\newblock \emph{Annals of Mathematics}, 181(1):155--190, 2015.

\bibitem{kerbl2023gaussian}
B.~Kerbl, G.~Kopanas, T.~Leimk\"uhler, and G.~Drettakis.
\newblock 3D Gaussian Splatting for real-time radiance field rendering.
\newblock \emph{ACM Transactions on Graphics (TOG)}, 42(4):1--14, 2023.

\end{thebibliography}

\end{document}
